% v3-figures/districts.tex: the 45 canonical tables of 01_schema.sql at 1.0.0-rc.1,
% arranged as rings around the Version 1 nucleus. The nucleus geometry (the
% hexagonal lattice of Version 1) is drawn faintly at the centre with the same
% proportions; the three outer rings are the tables the system grew. Every
% name is a real table. The ring numbers at the bottom match the caption.
\begin{tikzpicture}[scale=1.0,
  tbl/.style={draw=navymid, line width=0.6pt, rounded corners=1.5pt, fill=white, text=navydark,
              font=\tiny\ttfamily, inner sep=2.2pt, minimum width=2.5cm, minimum height=0.36cm, align=center},
  tblcore/.style={tbl, fill=navydark, text=white, font=\tiny\bfseries\sffamily},
  tblhub/.style={tbl, fill=accent, draw=accent, text=white, font=\scriptsize\bfseries\sffamily, minimum width=2.4cm, minimum height=0.55cm},
  tbladd/.style={tbl, dash pattern=on 1.6pt off 1.1pt, fill=tabstripe},
  ringnum/.style={font=\bfseries\scriptsize, text=gold, fill=white, draw=gold, line width=0.5pt, circle, inner sep=1.4pt},
]
\def\R{2.35}
\begin{scope}[on background layer]
  \foreach \a in {90,150,210,270,330,30} { \draw[steelgray, line width=0.4pt, opacity=0.10] (\a:\R) circle (0.62); }
  \draw[steelgray, line width=0.4pt, opacity=0.10] (0,0) circle (0.62);
  \foreach \a in {90,150,210,270,330,30} { \draw[steelgray, line width=0.3pt, opacity=0.07] (\a:{\R/2}) circle (0.42); }
  \draw[steelgray, line width=0.4pt, opacity=0.09] (90:\R) -- (150:\R) -- (210:\R) -- (270:\R) -- (330:\R) -- (30:\R) -- cycle;
  \draw[steelgray, line width=0.3pt, opacity=0.07] (90:{\R/2}) -- (150:{\R/2}) -- (210:{\R/2}) -- (270:{\R/2}) -- (330:{\R/2}) -- (30:{\R/2}) -- cycle;
  \draw[steelgray, line width=0.3pt, opacity=0.06] (90:\R) -- (210:\R) -- (330:\R) -- cycle;
  \draw[steelgray, line width=0.3pt, opacity=0.06] (270:\R) -- (30:\R) -- (150:\R) -- cycle;
  \draw[ring, opacity=0.55] (0,0) circle (5.2);
  \draw[ring, opacity=0.55] (0,0) circle (8.1);
  \draw[ring, opacity=0.55] (0,0) circle (11.0);
  \draw[ring, opacity=0.35, dash pattern=on 3pt off 2pt] (0,0) circle (13.9);
\end{scope}
% --- ring 1: the nucleus (Version 1's seven entities)
\node[tblhub] at (0,0) {ТокенЛичности};
\node[tblcore] at (90:\R)  {ФизическоеЛицо};
\node[tblcore] at (150:\R) {СобытиеЦиклаТокена};
\node[tblcore] at (30:\R)  {КонтекстПроверки};
\node[tblcore] at (210:\R) {Ведомство};
\node[tblcore] at (330:\R) {КриптоАлгоритм};
\node[tblcore] at (270:\R) {СобытиеПроверки};
% --- ring 2: the credential machinery (13)
\foreach \a/\n in {90/ПодписьТокена, 117.7/РазрешениеТокена, 145.4/ПривязкаУстройства, 173.1/СписокОтзывов, 200.8/ЗапросВосстановления,
                   228.5/СобытиеПринуждения, 256.2/СостояниеРегистрации, 283.8/ЭпохаТокенов, 311.5/ЛистЭпохиТокенов, 339.2/ЧислоДоказательства,
                   6.9/ЯкорьЦепиБлоков, 34.6/СобытиеКлючаДержателя, 62.3/ПерсонализацияКарты} {
  \pgfmathsetmacro{\rot}{ifthenelse(cos(\a)<-0.001,\a+180,\a)}
  \node[tbl, rotate=\rot] at (\a:5.2) {\n};
}
% --- ring 3: authority, trust, policy, and the people at the counter (13)
\foreach \a/\n in {90/ДопускКАлгоритму, 117.7/АттестацияДоверия, 145.4/СобытиеКлючаОргана, 173.1/ПолитикаУсмотрения, 200.8/КвотаВедомства,
                   228.5/ДоверяющаяСторона, 256.2/КодВходаПоглощён, 283.8/ЧислоОбмена, 311.5/ПользовательПриложения, 339.2/ПодтверждениеРегистрации,
                   6.9/ОснованияРегистрации, 34.6/Поручительство, 62.3/КодРегистрации} {
  \pgfmathsetmacro{\rot}{ifthenelse(cos(\a)<-0.001,\a+180,\a)}
  \node[tbl, rotate=\rot] at (\a:8.1) {\n};
}
% --- ring 4: archives, logs, and decisions kept as records. Twelve canonical tables and
%     the migration registry, which is the record of every schema change, so that this
%     ring has the same thirteen spokes as rings 2 and 3 and every spoke runs straight out.
\foreach \a/\n in {117.7/ПакетЯкорей, 145.4/ЖурналРасписок, 173.1/ЖурналОтметокВремени, 200.8/ЖурналВходов,
                   228.5/СобытиеСтиранияЛица, 256.2/КонтрольнаяТочкаАрхива, 283.8/ПолитикаХранения,
                   311.5/ПакетРегистраций, 339.2/БуферРегистраций, 6.9/СобытиеВедомства,
                   34.6/СобытиеПользователя, 62.3/СобытиеСтороны} {
  \pgfmathsetmacro{\rot}{ifthenelse(cos(\a)<-0.001,\a+180,\a)}
  \node[tbl, rotate=\rot] at (\a:11.0) {\n};
}
\node[tbladd, rotate=90] at (90:11.0) {ВерсияСхемы};
% --- ring 5: the rest of what a migrated deployment adds around the canonical schema,
%     six tables on a hexagon that echoes the nucleus, turned so the bottom axis stays free
\foreach \a/\n in {0/СессияОператора, 60/АппаратныйКлючОператора, 120/ЖурналДоступаАудита,
                   180/афина\_основное\_правило, 240/афина\_хранение\_ключей, 300/афина\_исполнение\_правил} {
  \pgfmathsetmacro{\rot}{ifthenelse(cos(\a)<-0.001,\a+180,\a)}
  \node[tbladd, rotate=\rot] at (\a:13.9) {\n};
}
% --- ring numbers on the bottom axis, in the gaps between boxes
\node[ringnum] at (0,-3.15) {1};
\node[ringnum] at (0,-5.2) {2};
\node[ringnum] at (0,-8.1) {3};
\node[ringnum] at (0,-11.0) {4};
\node[ringnum] at (0,-13.9) {5};
\end{tikzpicture}
